% !TEX TS-program = xelatex
% !TEX encoding = UTF-8 Unicode
% !Mode:: "TeX:UTF-8"

\documentclass{resume}
\usepackage{zh_CN-Adobefonts_external} % Simplified Chinese Support using external fonts (./fonts/zh_CN-Adobe/)
%\usepackage{zh_CN-Adobefonts_internal} % Simplified Chinese Support using system fonts
\usepackage{linespacing_fix} % disable extra space before next section
\usepackage{cite}
\usepackage[hidelinks]{hyperref}
\urlstyle{same}

\begin{document}
\pagenumbering{gobble} % suppress displaying page number

\name{陈坤}

\contactInfo{13018204668}{chenkun2024@ia.ac.cn}{}{https://github.com/kwen-chen}

\section{教育背景}
\datedsubsection{\textbf{电子科技大学}, 自动化, \textit{本科}, \textbf{GPA:} 3.94}{2020.09 - 2024.06}
\datedsubsection{\textbf{中国科学院自动化研究所}, 计算机应用技术, \textit{博士研究生（直博）}}{2024.09 - /}

\section{技术能力}
\begin{itemize}[parsep=0.2ex]
  \item \textbf{大模型、多模态大模型训练与对齐}: Python/PyTorch、继续预训练、指令微调、蒸馏、RL（Verl/Trl/OpenRLHF）、 Agentic RL、RAG、DeepResearch、数据自动合成（Self-Instruct/Self-QA）
  \item \textbf{并行与训练加速}: DP/MP/PP，Ray，DeepSpeed，FlashAttention，Ring-Attention，Ulysses
\end{itemize}

\section{实习经历}
\datedsubsection{\textbf{美团}, 基础算法研发部门算法实习生}{2025.08 - 2026.02}
\begin{itemize}
  \item 参与多模态大模型通用推理优化的研究：发表包括 RL冷启动泛化、RL梯度裁剪的两篇会议论文。
\end{itemize}

\datedsubsection{\textbf{中科闻歌}, AI研究院算法工程师}{2024.06 - 2025.06}
\begin{itemize}
  \item 长文本扩展（4K$\rightarrow$256K）与搜索智能体研发，支撑长文档场景落地；相关技术申请发明专利两项。
\end{itemize}

\section{学术成果}
\textbf{论文}（4篇： 2篇已录用，2篇在投，3篇为第一作者）
\begin{itemize}[parsep=0.2ex]
    \item \textbf{SPECS: Decoupling Multimodal Learning via Self-distilled Preference-based Cold Start}. \\\textit{ICLR 2026 Poster, Accepted}. 第一作者。arXiv:2510.25801（\href{https://arxiv.org/abs/2510.25801}{link}）
    \item \textbf{DEMO: Reframing Dialogue Interaction with Fine-grained Element Modeling}. \\\textit{ACL 2025 Findings, Accepted}. 第三作者（\href{https://aclanthology.org/2025.findings-acl.594/}{link}）
    \item \textbf{Flexible Entropy Control in RLVR with Gradient-Preserving Perspective}. \\\textit{ICML 2026, Under Review}. 第一作者。arXiv:2602.09782（\href{https://arxiv.org/abs/2602.09782}{link}）
    \item \textbf{APEX-Searcher: Augmenting LLMs' Search Capabilities through Agentic Planning and Execution}. \\\textit{SIGIR 2026, Under Review}. 第一作者
\end{itemize}

\textbf{专利}（4项：2项发明专利申请，2项已授权实用新型）
\begin{itemize}[parsep=0.2ex]
    \item \textbf{基于大模型任务规划的电力复杂任务执行方法及装置}. 发明专利申请，申请号：2025119951327
    \item \textbf{基于大语言模型的自动化长文本微调指令集构建方法}. 发明专利申请，申请号：202510119646X
    \item \textbf{一种移动式增材制造装置}. 实用新型（已授权），专利号：CN219988482U
    \item \textbf{一种增减材一体式制造设备}. 实用新型（已授权），专利号：CN220146695U
\end{itemize}

\section{科研项目经历}
\datedsubsection{\textbf{多模态出版内容资源编目与融合检索技术研究（2024JK4003）}}{}
\begin{itemize}
    \item 参与多模态出版内容资源加工平台研发，支撑多源异构资源的智能化全流程处理。
    \item 完成相关大模型能力训练、自我认知训练与能力评测工作。
\end{itemize}

\datedsubsection{\textbf{国家电网科技项目：大语言模型关键技术研究及在电力智能客服中的示范应用}}{}
\begin{itemize}
    \item 参与大模型训练、数据构建与技术文档撰写，并作为主要执笔人申请发明专利两项。
\end{itemize}

\datedsubsection{\textbf{中科闻歌大语言模型长上下文扩展技术研究项目}}{}
\begin{itemize}
    \item 基于 DeepSpeed Ulysses 与 Ring Attention 等序列并行方案，将训练上下文从4K扩展至256K，并开源框架 Eazycontext（\textbf{731 stars}）。
\end{itemize}

\section{荣誉/获奖经历}
\begin{itemize}[parsep=0.2ex]
    \item \textbf{电子科技大学成电杰出学生} (\textit{2023.12}) 全校\textbf{10名}本科生入选
    \item \textbf{美国大学生数学建模竞赛 Meritorious Winner（美赛一等奖）} (\textit{2023.05}) 
    \item \textbf{第十八届"挑战杯"全国大学生课外学术科技作品科技发明制作专项国家级一等奖} (\textit{2023.10})
    \item \textbf{第十八届全国大学生智能车竞赛国家级一等奖} (\textit{2023.8})
    \item \textbf{多次获得国家奖学金、深交所奖学金等} (\textit{2020-2025})
\end{itemize}

% This section is intentionally left blank.

%% Reference
%\newpage
%\bibliographystyle{IEEETran}
%\bibliography{mycite}
\end{document}
